% ======================================================
%  The one page that is this volume's alone: what it assumes.
%
%  It exists because the honest answer to "which book do I buy" is not
%  "the harder one", and a reader who starts here and bounces off will
%  conclude the method does not work rather than that they opened the
%  wrong volume.
% ======================================================

\chapter*{Dla kogo jest ta książka}
\addcontentsline{toc}{chapter}{Dla kogo jest ta książka}
\markboth{Dla kogo jest ta książka}{Dla kogo jest ta książka}

To jest \textbf{drugi tom}. Pierwszy — \emph{Trening Mózgu} — ćwiczy rachunki:
dodawanie, mnożenie, procenty, ułamki dziesiętne, kolejność działań. Ten nie
ćwiczy ich jeszcze raz na większych liczbach.

\section*{Czego tu nie ma}

\textbf{Nie ma dłuższych sum.} Pierwszy tom kończy się na dodawaniu
pięciocyfrowym, mnożeniu trzycyfrowego przez dwucyfrowe i procencie składanym.
Dopisanie jeszcze jednej cyfry daje działanie \emph{dłuższe}, a nie
trudniejsze — i dziewięćdziesiąt zestawów tego byłoby tą samą książką
wydrukowaną dwa razy.

\section*{Czego tu jest}

Działania, których w pierwszym tomie nie ma wcale:

\begin{center}
  \begin{tabular}{@{}ll@{}}
    \textbf{Systemy i reszty} & dwójkowy, szesnastkowy, modulo, cyfra kontrolna \\[3pt]
    \textbf{Struktura liczby} & czynniki pierwsze, dzielniki, Euklides \\[3pt]
    \textbf{Ułamki i potęgi} & mnożenie i dzielenie ułamków, wykładnik ujemny \\
                             & i ułamkowy, logarytmy, notacja naukowa \\[3pt]
    \textbf{Zastosowania}    & kombinatoryka, dzień tygodnia z daty, prędkość \\
                             & średnia, skala mapy, procent przez lata \\
  \end{tabular}
\end{center}

\noindent
Reguły gry są te same: czterdzieści zadań w zestawie, w pamięci, na stoper,
wszystkie odpowiedzi na końcu książki.

\section*{Skąd wiadomo, że pora}

Nie ma progu i nie ma egzaminu. Praktyczny test jest taki: \textbf{jeśli
zestawy z Bloku IV pierwszego tomu mieścisz w czasie docelowym i nie musisz się
przy nich zatrzymywać}, ten tom nie będzie za wcześnie.

A jeśli będzie — to nic nie kosztuje. Odłóż go i wróć za miesiąc; pierwszy tom
ma materiału na cztery.
